\documentclass[11pt,a4paper]{article}
\usepackage[margin=1.55cm,top=1.35cm,bottom=1.35cm]{geometry}
\usepackage{amsmath,amssymb,mathtools}
\usepackage{xcolor}
\usepackage{array}
\usepackage{enumitem}
\usepackage{microtype}

\setlength{\parindent}{0pt}
\setlength{\parskip}{5pt}
\setlist[enumerate]{leftmargin=1.2cm,itemsep=2pt,topsep=2pt}
\allowdisplaybreaks

\definecolor{noteRed}{RGB}{210,25,25}
\newcommand{\rednote}[1]{\textcolor{noteRed}{#1}}
\newcommand{\D}{\mathrm{D}}
\newcommand{\C}{\mathrm{C}}

\begin{document}
\pagestyle{empty}

% Source page 2
\underline{\textbf{Subgroups:}}

\underline{Definition:}\quad $H$ of $G$ is a subset of $G$ satisfying the group axioms

Two trivial subgroups: $\{e\}$ and $G$

If $g\in G$, the smallest subgroup containing $g$ is
\[
\langle g\rangle=\{g^n\mid n\in\mathbb{Z}\}.
\]

\emph{Ex:} $\D_3$
\[
\{e\},\qquad \D_3,\qquad \{e,r,r^2=r^{-1}\},\qquad \{e,\sigma_1\}.
\]

\underline{\textbf{Normal subgroup:}}

A subgroup $H$ is normal (invariant)
if
\[
gHg^{-1}=H,\qquad \forall g\in G.
\]
\[
\boxed{H\triangleleft G}
\]

\emph{Ex:} In $\D_3$: $\{e,r,r^{-1}\}$ is the only proper normal subgroup.

\underline{\textbf{Simple and Semi-simple groups:}}

\underline{Simple:}\quad No nontrivial normal subgroups

\underline{Semi-Simple:}\quad No abelian normal subgroups

\underline{\textbf{Coset decomposition and Lagrange's theorem}}

$H$ is a subgroup of $G$. Cosets are a way of partitioning the group into disjoint pieces of equal size, each piece being ``a shifted copy'' of the subgroup.

% Source page 3
For $\forall a\in G$,
\[
Ha=\{ha\mid h\in H\},\qquad \text{right coset}
\]
\[
aH=\{ah\mid h\in H\},\qquad \text{left coset}
\]

\underline{Property:}\quad $Ha_i=Ha_j$ or $Ha_i\cap Ha_j=\varnothing$

\underline{Property:}\quad If $H$ is normal,
\[
aHa^{-1}=H,
\]
\[
aH=Ha,
\]
all left and right coset coincide \ldots

\begin{center}
\underline{\textbf{Lagrange theorem}}
\end{center}

Let $G$ be a finite group, $H$ a subgroup. Then:

\begin{enumerate}
\item
\[
G=Ha_1\cup Ha_2\cup\cdots\cup Ha_n,\qquad a_1=e,
\]
\[
G=b_1H\cup b_2H\cup\cdots\cup b_nH,\qquad b_1=e.
\]

\item
\[
|Ha_i|=|a_iH|=|H|.
\]

\item The number of cosets is the \underline{index} of $H$ in $G$:
\[
|G|=n|H|,
\qquad
n=|G:H|.
\]
\end{enumerate}

\emph{Ex:}
\[
\D_3=\{e,r,r^{-1},\sigma_1,\sigma_2,\sigma_3\},\qquad |\D_3|=6,
\]
\[
H=\{e,\sigma_1\},\qquad |H|=2
\]
\[
\Longrightarrow n=|G:H|=\frac{6}{2}=3.
\]

% Source page 4
\underline{\textbf{Let's compute left cosets}}
\[
eH=\{e,\sigma_1\}=H
\]
\[
rH=\{re,r\sigma_1\}=\{r,\sigma_3\}
\]
\[
r^{-1}H=\{r^{-1},r^{-1}\sigma_1\}=\{r^{-1},\sigma_2\}
\]
\[
\sigma_2H=\{\sigma_2,\sigma_2\sigma_1\}=r^{-1}H
\]
\[
\sigma_3H=rH
\]
\[
\sigma_1H=eH
\]

\[
\D_3=\{e,\sigma_1\}\cup\{r,\sigma_3\}\cup\{r^{-1},\sigma_2\}
\]
\[
\Longrightarrow\quad\text{for a right coset:}
\]
\[
\D_3=\{e,\sigma_1\}\cup\{r,\sigma_2\}\cup\{r^{-1},\sigma_3\}.
\]

\underline{\textbf{Exercise}}
\[
K=\{e,r,r^{-1}\},
\qquad
gKg^{-1}=K,\quad\forall g\in\D_3
\]
\[
gK=Kg
\]
\[
|K|=3\quad\Longrightarrow\quad n=\frac{6}{3}=2.
\]

\[
srs=r^{-1},\qquad srs^{-1}=r^{-1},\qquad r^3=e.
\]

\[
eK=K
\]
\[
rK=\{re,rr,rr^{-1}\}=\{r,r^2,e\}=K
\]
\[
r^2K=\{r^2e,r^2r,r^2r^{-1}\}=\{e,r,r^2\}=K
\]
\[
\sigma_1K=\{\sigma_1e,\sigma_1r,\sigma_1r^2\}
=\{\sigma_1,\sigma_3,\sigma_2\}
\]
\[
\cdots
\]
\[
\sigma_2K=\sigma_3K=\sigma_1K=\{\sigma_1,\sigma_2,\sigma_3\}
\]
\[
\D_3=K\cup\sigma_1K
\qquad
\underset{\text{rotations}}{K}\qquad
\underset{\text{reflections}}{\sigma_1K}.
\]

% Source page 5
\underline{\textbf{Conjugate class decomposition}}

\underline{\textbf{Conjugate elements:}}\quad $a,b\in G$; $b$ is conjugate of $a$ if there exists $g\in G$ such that
\[
b=gag^{-1}.
\]

\underline{Prop:}
\begin{enumerate}
\item $a=eae^{-1}$.

\item If $b=gag^{-1}$,
\[
a=g^{-1}bg.
\]

\item If $b=g_1ag_1^{-1}$, $c=g_2bg_2^{-1}$,
\[
c=g_2g_1a\,(g_2g_1)^{-1}.
\]
\end{enumerate}

\underline{\textbf{Conjugacy classes:}}

The conjugacy class of an element $a\in G$ is the set
\[
(a)=\{gag^{-1}\mid g\in G\}.
\]

\rednote{Prop:}
\begin{enumerate}
\item If $(a)=\{a\}$, then $a$ is self-conjugate.

\item In any abelian groups, every element is self-conjugate
\[
[gag^{-1}=a].
\]

\item Conjugacy classes are invariant under conjugation:
\[
g(a)g^{-1}=(a),\qquad\forall g\in G.
\]
\end{enumerate}

\emph{Ex:} $\D_3$
\[
(e)=\{e\}\quad\Longleftrightarrow\quad\text{exercise}
\]
\[
(r)=(r^{-1})=\{r,r^{-1}\}
\]
\[
(\sigma_1)=(\sigma_2)=(\sigma_3)=\{\sigma_1,\sigma_2,\sigma_3\}
\]

\underline{\textbf{Theorem:}} Every group $G$ is the disjoint union of the conjugacy classes of its elements
\[
G=\bigcup_{a\in G}(a)
\]
\[
\D_3=\{e\}\cup\{r,r^{-1}\}\cup\{\sigma_1,\sigma_2,\sigma_3\}.
\]

% Source page 6
\underline{\textbf{Quotient group:}}\quad $H$ is a normal subgroup of $G$ ($H\triangleleft G$)

Assume
\[
G=a_1H\cup a_2H\cup\cdots\cup a_nH,\qquad a_1=e.
\]

The set
\[
G/H=\{a_1H,a_2H,\ldots,a_nH\}
\]
is called the quotient group of $G$ by $H$.

\underline{Prop:}
\begin{enumerate}
\item For $a_iH,a_jH\in G/H$ the composition is defined as
\[
(a_iH)(a_jH)=(a_ia_j)H.
\]

\item $eH=He=H$.

\item
\[
(a_iH)^{-1}=a_i^{-1}H.
\]

\item $|G|=n_G$, $|H|=n_H$
\[
\Longrightarrow\qquad |G/H|=\frac{n_G}{n_H}.
\]
\end{enumerate}

\emph{Ex:} $\D_3$
\[
H=\{e,r,r^{-1}\}
\]
\[
eH=rH=r^{-1}H
\]
\[
\sigma_1H=\sigma_2H=\sigma_3H=\{\sigma_1,\sigma_2,\sigma_3\}
\]
\[
\D_3/H=\{H,\sigma_1H\}.
\]

\[
\begin{array}{c|cc}
\D_3/H & H & \sigma_1H\\ \hline
H & (H)(H)=H & (H)(\sigma_1H)=\sigma_1H\\[2pt]
\sigma_1H & (\sigma_1H)(H)=\sigma_1H &
(\sigma_1H)(\sigma_1H)=\sigma_1\sigma_1H=H
\end{array}
\]

\[
C_2=\{e,g\},\qquad g^2=e,
\qquad \D_3/H\simeq C_2.
\]

% Source page 7
\underline{\textbf{Centraliser:}}

The centraliser of $S$ in $G$ is the set of
\[
C_G(S)=\{g\in G\mid gs=sg\text{ for all }s\in S\}.
\]

\underline{\textbf{Normalizer:}}
\[
N_G(S)=\{g\in G\mid gSg^{-1}=S\}
=\{g\in G\mid gsg^{-1}\in S\text{ for all }s\in S\}.
\]

\underline{\textbf{Center:}}

Center of $G$:
\[
Z(G)=\{z\in G\mid zg=gz\text{ for all }g\in G\}.
\]

\underline{\textbf{Remarks}}
\begin{enumerate}
\item $C_G(S)$ and $N_G(S)$ are subgroups of $G$.

\item $Z(G)$ is always abelian subgroup of $G$.

\item If $S$ is itself a subgroup then $C_G(S)$ measures how close $S$ is to being central -- whereas $N_G(S)$: how close $S$ is to being normal.
\end{enumerate}

\underline{\textbf{Direct product}}

A group $G$ is the direct product of $H$ and $K$, $G=H\times K$ if:
\begin{enumerate}
\item Every $h\in H$ commutes with $k\in K$, $hk=kh$.

\item Every $g\in G$ can be written \underline{uniquely} as $g=hk$.
\end{enumerate}

\underline{Prop.}
\[
|H|=n_H,\qquad |K|=n_K
\]
\[
\Longrightarrow |G|=n_H\times n_K.
\]

\underline{Composition}
\[
(h_1,k_1)(h_2,k_2)=(h_1h_2,k_1k_2).
\]

\emph{Ex:}
\[
G=\D_3\times\mathbb{Z}_2^+,\qquad
\mathbb{Z}_2^+=\{0,1\},\quad\text{addition}
\]
\[
\Longrightarrow\quad\text{exercise: Show }\D_6=\D_3\times\mathbb{Z}_2^+.
\]

\end{document}
